\section{Lecture 1: Simple Linear Regression}

\subsection{Regression: Basic Concepts}

Regression aims to find the \emph{best} function
\[
f : \mathbb{R}^p \to \mathbb{R}
\]
such that
\[
Y \approx f(x_1, \dots, x_p), \qquad p \ge 1.
\]

\begin{itemize}
    \item $Y$ is called the \textbf{response} or \textbf{dependent variable}.
    \item $x_1, \dots, x_p$ are called \textbf{explanatory}, \textbf{independent}, or \textbf{predictor variables}.
    \item Predictors are also referred to as \textbf{covariates} or \textbf{features}.
\end{itemize}

\subsection{Simple vs.\ Multiple Regression}

\begin{itemize}
    \item When $p = 1$:
    \[
    Y \approx f(x) \quad \Rightarrow \quad \textbf{Simple Regression}.
    \]
    \item When $p > 1$:
    \[
    Y \approx f(x_1, \dots, x_p) \quad \Rightarrow \quad \textbf{Multiple Regression}.
    \]
\end{itemize}

In this lecture, we focus on \textbf{simple regression}.

\subsection{Parametric Regression}

In general, it is difficult to determine the ``best'' function $f$.

To make the problem tractable, we \emph{specify the form of $f$ using a finite number of parameters}.
This approach is called \textbf{parametric regression}.

\subsection{Examples}

\begin{enumerate}
    \item Constant model:
    \[
    f(x) = \beta_0
    \]
    \item Linear model:
    \[
    f(x) = \beta_0 + \beta_1 x
    \]
    \item Quadratic model:
    \[
    f(x) = \beta_0 + \beta_1 x + \beta_2 x^2
    \]
\end{enumerate}

In each case, the function $f$ depends on unknown parameters $\beta_0, \beta_1, \beta_2, \ldots$

\subsection{Linear Part of the Model}

On average, the response is linear in the parameters:
\[
\beta_0 + \beta_1 x_{i1} + \cdots + \beta_p x_{ip}.
\]

This expression is linear in $\boldsymbol{\beta}$. However, the entries of the predictor vector $\mathbf{x}$ may include:
\begin{itemize}
    \item Continuous or discrete predictors.
    \item Dummy variables for categorical predictors.
    \item Polynomial terms such as $x^2$, $x^3$, etc.
    \item Transformations such as $\log(x)$ or $\sin(x)$.
    \item Interaction or cross-product terms such as $x_1 x_2$.
\end{itemize}

\subsection{Simple Linear Regression Model}

Given observed data
\[
(x_1, Y_1), \dots, (x_n, Y_n),
\]
we propose the model
\[
Y_i = f(x_i) + \epsilon_i,
\]
where
\[
f(x_i) = \beta_0 + \beta_1 x_i
\quad \text{and} \quad
\E[\varepsilon_i] = 0, \qquad i = 1, \dots, n.
\]

\subsection{Model Components}

\begin{itemize}
    \item \textbf{Intercept parameter} $\beta_0$: deterministic and unknown
    \item \textbf{Slope parameter} $\beta_1$: deterministic and unknown
    \item \textbf{Error term} $\varepsilon_i$: random, zero-mean, and unknown
\end{itemize}

The intercept and slope are \emph{common across all observations}, while the errors vary between data points.

\subsection{Interpretation Diagram (Conceptual)}

\[
Y_i = \beta_0 + \beta_1 x_i + \varepsilon_i
\]

\begin{itemize}
    \item $Y_i$: response variable
    \item $x_i$: explanatory variable
    \item $\beta_0$: intercept
    \item $\beta_1$: slope
    \item $\varepsilon_i$: error
\end{itemize}

\subsection{Formal Properties of the Simple Linear Regression Model}

The simple linear regression model is
\[
Y_i = \beta_0 + \beta_1 x_i + \varepsilon_i, \qquad i = 1, \dots, n.
\]

\begin{itemize}
    \item The designation \textbf{simple} indicates there is only one predictor variable $x$.
    \item The model is called \textbf{linear} because it is linear in the parameters $\beta_0$ and $\beta_1$.
\end{itemize}

\subsection{Expectation}

Formally, given $X = x_i$, we have
\[
\begin{aligned}
\E(Y_i)
&= \E(\beta_0 + \beta_1 x_i + \varepsilon_i) \\
&= \E(\beta_0 + \beta_1 x_i) + \E(\varepsilon_i) \\
&= \beta_0 + \beta_1 x_i,
\end{aligned}
\]
since $\E(\varepsilon_i) = 0$.

Thus,
\[
\E(Y_i \mid X = x_i) = \beta_0 + \beta_1 x_i.
\]

\subsection{Linearity in Regression Models}

Consider the following regression models.

\subsubsection*{Examples of Nonlinear Models (in Parameters)}

\[
Y = \exp(\beta_0 + \beta_1 x) + \varepsilon
\]

This model is \textbf{not linear} in $\beta_0$ and $\beta_1$.

\[
Y = \beta_0 + \beta_1^2 x + \varepsilon
\]

This model is also \textbf{not linear} in $\beta_0$ and $\beta_1$ due to the squared parameter.

\subsubsection*{Examples of Linear Models (in Parameters)}

\[
Y = \beta_0 \cdot 1 + \beta_1 x + \beta_2 \log(x) + \varepsilon
\]

This model is \textbf{linear} in the parameters $\beta_0$, $\beta_1$, and $\beta_2$.

\[
Y = \beta_0 + \beta_1 \log(x) + \varepsilon
\]

This is a \textbf{simple linear regression model}, since there is only one predictor $\log(x)$ and the model is linear in the parameters.

\subsection{Sample Level vs.\ Population Level}

\subsubsection*{Sample Level}

The observed data
\[
(x_1, Y_1), (x_2, Y_2), \dots, (x_n, Y_n)
\]
are realizations of the random variable $(X, Y)$.

\subsubsection*{Population Level}

At the population level, the random variable $(X, Y)$ follows a joint distribution satisfying
\[
Y = \beta_0 + \beta_1 X + \varepsilon.
\]

\subsection{Modeling Assumption}

In this course, we treat $X$ as given (known) and $Y$ as a random variable for simplicity.

Under this treatment,
\[
\E(Y \mid X) = \E(Y),
\qquad
\Var(Y \mid X) = \Var(Y).
\]

The randomness of $Y_i$ arises from the error term $\varepsilon_i$, since
\[
Y_i = \beta_0 + \beta_1 x_i + \varepsilon_i.
\]

\subsection{Fitted Regression Model}

For simple linear regression, the \textbf{fitted regression model} is
\[
\hat{Y} = \bhat{\beta}_0 + \bhat{\beta}_1 x,
\]
where $\bhat{\beta}_0$ and $\bhat{\beta}_1$ are estimates of the unknown parameters $\beta_0$ and $\beta_1$, respectively.

\subsection{Interpretation}

The fitted value $\hat{Y}$ estimates the conditional mean of the response:
\[
\hat{Y} \approx \E(Y \mid X).
\]

\subsection{Clarification}

\[
Y = \beta_0 + \beta_1 x + \varepsilon
\qquad \text{(true model)}
\]

\[
\E(Y \mid X) = \beta_0 + \beta_1 x
\qquad \text{(population mean)}
\]

\[
\hat{Y} = \bhat{\beta}_0 + \bhat{\beta}_1 x
\qquad \text{(estimated mean)}
\]

\subsection{Interpretation of Regression Coefficients (Pearson's Data)}

Consider the fitted regression model
\[
\widehat{\text{son}} = 33.89 + 0.51 \times \text{father},
\]
where
\begin{itemize}
    \item \textbf{son} = height of sons (in inches),
    \item \textbf{father} = height of fathers (in inches).
\end{itemize}

\subsubsection*{Intercept}

When $x = 0$, the expected value of the response is $33.89$.

\textbf{Question:} Does it make sense to assume that a father's height is $0$ inches in this context?

\textbf{Interpretation:}
On average, a son's height is $33.89$ inches when the father's height is zero (not meaningful in practice).

\subsubsection*{Slope}

When $x$ increases by 1 unit, the expected value of the response increases by $0.51$ on average.

\textbf{Contextual Meaning:}
In Pearson's data, one unit corresponds to \textbf{one inch of height}.

\textbf{Interpretation:}
When a father's height increases by 1 inch, the son's height increases by $0.51$ inches on average.

\subsection{Useful Algebraic Identities}

We frequently use the following identities.

\subsubsection*{Centered Cross-Product}

\[
\sum_{i=1}^n (x_i - \bar{x})(y_i - \bar{y})
=
\sum_{i=1}^n x_i y_i - n \bar{x}\bar{y}.
\]

\subsubsection*{Centered Sum of Squares}

\[
\sum_{i=1}^n (x_i - \bar{x})^2
=
\sum_{i=1}^n x_i^2 - n \bar{x}^2,
\]

\[
\sum_{i=1}^n (y_i - \bar{y})^2
=
\sum_{i=1}^n y_i^2 - n \bar{y}^2.
\]

\subsubsection*{Sample Means}

\[
\bar{x} = \frac{1}{n}\sum_{i=1}^n x_i
\quad \Rightarrow \quad
\sum_{i=1}^n x_i = n\bar{x},
\]

\[
\bar{y} = \frac{1}{n}\sum_{i=1}^n y_i
\quad \Rightarrow \quad
\sum_{i=1}^n y_i = n\bar{y}.
\]

\subsection{Least Squares Estimation}

Recall the simple linear regression model:
\[
Y_i = \beta_0 + \beta_1 x_i + \varepsilon_i.
\]

\begin{defn}{Method of Least Squares}{def:leastsquares}
The \textbf{method of least squares} estimates $\beta_0$ and $\beta_1$ by minimizing the sum of squared residuals:
\[
\min_{\beta_0,\beta_1} L(\beta_0,\beta_1),
\]
where
\[
L(\beta_0,\beta_1) = \sum_{i=1}^n (y_i - \beta_0 - \beta_1 x_i)^2.
\]
\end{defn}

Minimizing $L(\beta_0,\beta_1)$ means finding the values of $\beta_0$ and $\beta_1$ that make this function as small as possible.

\subsection{Normal Equations}

\subsubsection*{Step 1: Partial Derivatives}

Take derivatives of $L(\beta_0,\beta_1)$ with respect to $\beta_0$ and $\beta_1$:

\[
\frac{\partial L}{\partial \beta_0}
=
-2\sum_{i=1}^n (y_i - \beta_0 - \beta_1 x_i),
\]

\[
\frac{\partial L}{\partial \beta_1}
=
-2\sum_{i=1}^n x_i (y_i - \beta_0 - \beta_1 x_i).
\]

\subsubsection*{Step 2: Set Derivatives Equal to Zero}

\[
\sum_{i=1}^n (y_i - \beta_0 - \beta_1 x_i) = 0,
\]

\[
\sum_{i=1}^n x_i (y_i - \beta_0 - \beta_1 x_i) = 0.
\]

\subsection{Solving for the Intercept}

From
\[
\sum_{i=1}^n (y_i - \beta_0 - \beta_1 x_i) = 0,
\]
we obtain
\[
\sum_{i=1}^n y_i - n\beta_0 - \beta_1 \sum_{i=1}^n x_i = 0.
\]

Using $\sum y_i = n\bar{y}$ and $\sum x_i = n\bar{x}$:
\[
n\bar{y} - n\beta_0 - n\beta_1 \bar{x} = 0.
\]

Dividing by $n$:
\[
\beta_0 = \bar{y} - \beta_1 \bar{x}.
\]

Thus, the least squares estimator of the intercept is
\[
\bhat{\beta}_0 = \bar{y} - \bhat{\beta}_1 \bar{x}.
\]

If $\bhat{\beta}_1 = 0$, then $\bhat{\beta}_0 = \bar{y}$.

\subsection{Solving for the Slope}

Starting from
\[
\sum_{i=1}^n x_i (y_i - \beta_0 - \beta_1 x_i) = 0,
\]
substitute $\beta_0 = \bar{y} - \beta_1 \bar{x}$.

After simplification, we obtain
\[
\bhat{\beta}_1
=
\frac{\sum_{i=1}^n (x_i - \bar{x})(y_i - \bar{y})}
     {\sum_{i=1}^n (x_i - \bar{x})^2}.
\]

\subsubsection*{Interpretation}

\begin{itemize}
    \item Numerator: sample covariance of $x$ and $y$
    \item Denominator: sample variance of $x$
\end{itemize}

The solution depends entirely on the observed data.

\begin{keyresult}[Least Squares Estimators for Simple Linear Regression]
\[
\bhat{\beta}_1
=
\frac{\sum_{i=1}^n (x_i - \bar{x})(y_i - \bar{y})}
     {\sum_{i=1}^n (x_i - \bar{x})^2},
\]

\[
\bhat{\beta}_0
=
\bar{y} - \bhat{\beta}_1 \bar{x}.
\]
\end{keyresult}

\subsection{Residuals and Residual Sum of Squares}

\subsubsection*{Fitted Values}

For each observation $i = 1, \dots, n$, the fitted (predicted) value is
\[
\hat{y}_i = \bhat{\beta}_0 + \bhat{\beta}_1 x_i.
\]

\subsubsection*{Residuals}

The residual is defined as the difference between the observed and fitted values:
\[
\hat{e}_i = y_i - \hat{y}_i, \qquad i = 1, \dots, n.
\]

Residuals measure how far each observation lies from the fitted regression line.

\subsubsection*{Residual Sum of Squares}

\begin{keyresult}[Residual Sum of Squares]
The \textbf{Residual Sum of Squares} is defined by
\[
\RSS
=
\sum_{i=1}^n \hat{e}_i^2
=
\sum_{i=1}^n (y_i - \hat{y}_i)^2.
\]
\end{keyresult}

\begin{notebox}
\textbf{Remarks about RSS:}
\begin{itemize}
    \item RSS is differentiable and has convenient mathematical properties.
    \item Squaring penalizes large residuals more heavily than small ones.
\end{itemize}
\end{notebox}

\subsubsection*{Residual Vector Representation}

For observations $i = 1, \dots, n$:
\[
\hat{e}_1 = y_1 - \hat{y}_1,
\quad
\hat{e}_2 = y_2 - \hat{y}_2,
\quad \dots \quad
\hat{e}_n = y_n - \hat{y}_n.
\]

Thus,
\[
\RSS
=
\hat{e}_1^2 + \hat{e}_2^2 + \cdots + \hat{e}_n^2.
\]

\subsection{Properties of the Least Squares Fit}

\subsubsection*{Residuals Sum to Zero}

For the least squares estimators,
\[
\sum_{i=1}^n \hat{e}_i
=
\sum_{i=1}^n \bigl(y_i - (\bhat{\beta}_0 + \bhat{\beta}_1 x_i)\bigr)
=
0.
\]

\textbf{Proof:}

\[
\sum_{i=1}^n \hat{e}_i
=
\sum_{i=1}^n y_i
-
\sum_{i=1}^n \bhat{\beta}_0
-
\sum_{i=1}^n \bhat{\beta}_1 x_i.
\]

Using
\[
\sum_{i=1}^n y_i = n\bar{y},
\quad
\sum_{i=1}^n \bhat{\beta}_0 = n\bhat{\beta}_0,
\quad
\sum_{i=1}^n x_i = n\bar{x},
\]
we obtain
\[
n\bar{y} - n\bhat{\beta}_0 - n\bhat{\beta}_1 \bar{x}.
\]

Substituting $\bhat{\beta}_0 = \bar{y} - \bhat{\beta}_1 \bar{x}$:
\[
n\bar{y} - n(\bar{y} - \bhat{\beta}_1 \bar{x}) - n\bhat{\beta}_1 \bar{x} = 0.
\]

Hence,
\[
\sum_{i=1}^n \hat{e}_i = 0.
\]

\subsubsection*{Residuals Weighted by the Regressor}

Another fundamental property is
\[
\sum_{i=1}^n x_i \hat{e}_i = 0.
\]

\textbf{Interpretation:}

\begin{itemize}
    \item Residuals are orthogonal to the regressor $x$.
    \item This condition arises from the normal equations.
    \item It ensures the slope estimator is optimal in the least squares sense.
\end{itemize}

\subsubsection*{Exercise}

Show that
\[
\sum_{i=1}^n x_i \hat{e}_i = 0
\]
using the normal equations for $\bhat{\beta}_0$ and $\bhat{\beta}_1$.

\subsubsection*{Solution}

Recall that the residuals are
\[
\hat{e}_i = y_i - (\bhat{\beta}_0 + \bhat{\beta}_1 x_i).
\]

Then
\[
\sum_{i=1}^n x_i \hat{e}_i
=
\sum_{i=1}^n x_i \bigl(y_i - \bhat{\beta}_0 - \bhat{\beta}_1 x_i\bigr).
\]

Expanding the sum:
\[
\sum_{i=1}^n x_i y_i
-
\bhat{\beta}_0 \sum_{i=1}^n x_i
-
\bhat{\beta}_1 \sum_{i=1}^n x_i^2.
\]

From the normal equations for least squares:
\[
\sum_{i=1}^n x_i (y_i - \bhat{\beta}_0 - \bhat{\beta}_1 x_i) = 0.
\]

This expression is exactly the same as $\sum_{i=1}^n x_i \hat{e}_i$. Therefore,
\[
\sum_{i=1}^n x_i \hat{e}_i = 0.
\]

Hence, the residuals are orthogonal to the regressor variable $x$.

\clearpage
